\section*{Chapitre \ref{ch:b3:field-lab-methods-3} --- Méthodes de terrain et de laboratoire III}
% Values from the Book 5 scratch calculator calc/ch31.py.

\begin{solution}{exo:b3:field-lab-methods-3:1}
(a)~$\qty{0.01}{\milli\gal} = \qty{e-7}{m.s^{-2}}$, et
$\num{e-7}/\num{9.81} = \num{1.0e-8}$~: une partie sur cent millions.

(b)~$0.01/0.3086 = \qty{0.032}{m}$, soit environ \qty{3.2}{cm}.

(c)~Ensemble, les deux effets modifient la pesanteur de
\qty{0.1967}{\milli\gal.m^{-1}}, soit la différence $0.3086 - 0.1119$, donc
la hauteur doit varier de $0.01/0.1967 = \qty{5.1}{cm}$.
\end{solution}

\begin{solution}{exo:b3:field-lab-methods-3:2}
$\Delta m = \tfrac43\pi(\qty{50}{m})^3\times\qty{500}{kg.m^{-3}} =
\qty{2.62e8}{kg}$. Le pic vaut $G\Delta m/z^2 =
\num{6.674e-11}\times\num{2.62e8}/\num{e4} = \qty{1.75e-6}{m.s^{-2}}$,
soit \qty{0.175}{\milli\gal}. La demi-largeur vaut $x_{1/2} = z/1.305 =
\qty{77}{m}$.
\end{solution}

\begin{solution}{exo:b3:field-lab-methods-3:3}
(a)~Avec $AM = MN = NB = a$, $K = 2\pi/(1/a - 1/2a - 1/2a + 1/a) =
2\pi a = \qty{62.83}{m}$, et $\rho_a = 62.83\times0.239/0.1 =
\qty{150}{\ohm\metre}$. (b)~$K = \pi(15^2 - 5^2)/(2\times5) = \pi\times
200/10 = \qty{62.83}{m}$. Les électrodes sont alors en $0$, $10$, $20$
et \qty{30}{m}, dans l'ordre AMNB~: c'est le dispositif Wenner
d'espacement \qty{10}{m}.
\end{solution}

\begin{solution}{exo:b3:field-lab-methods-3:4}
(a)~$V = 0.2998/\sqrt5 = \qty{0.134}{m.ns^{-1}}$. (b)~$t =
2\times10/0.134 = \qty{149}{ns}$. (c)~$\delta = 0.0053\sqrt5/\num{e-3} =
\qty{11.9}{m}$, et une profondeur utile d'environ \qty{36}{m}. (d)~Dans l'argile
$\delta = 0.0053\sqrt{20}/0.01 = \qty{2.4}{m}$, une profondeur utile
d'environ \qty{7}{m}, et une vitesse de seulement \qty{0.067}{m.ns^{-1}}.
Ici, l'approximation reste valable~: à \qty{100}{MHz},
$\omega\varepsilon = 2\pi\times\num{e8}\times20\times\num{8.85e-12} =
\qty{0.11}{S.m^{-1}}$, dix fois $\sigma$.
\end{solution}

\begin{solution}{exo:b3:field-lab-methods-3:5}
(a)~Le cylindre porte $\lambda = \pi R^2\Delta\rho$ par unité de longueur, et
à l'extérieur, son attraction est celle d'une masse linéique sur son axe (théorème de Gauss
appliqué à un cylindre coaxial)~: $\Delta g = 2G\pi R^2\Delta\rho\,z/(x^2
+ z^2)$. La moitié du pic est atteinte lorsque $x^2 + z^2 = 2z^2$, soit $z =
x_{1/2}$. (b)~La sphère est à $1.305\times100 = \qty{130}{m}$, le
cylindre à \qty{100}{m}~: le même profil, interprété comme une sphère, place le
corps \qty{30}{\percent} plus profond. L'anomalie d'une sphère décroît plus vite avec
la distance, si bien qu'une sphère doit être plus profonde pour être aussi étendue.
\end{solution}

\begin{solution}{exo:b3:field-lab-methods-3:6}
(a)~Chaque point représente $\qty{100}{m}\times\qty{100}{m} =
\qty{e4}{m^2}$, soit $\iint\Delta g\,\mathrm{d}A \approx
\num{-438e-5}\times\num{e4} = \qty{-43.8}{m^3.s^{-2}}$, et $\Delta m
= -43.8/(2\pi\times\num{6.674e-11}) = \qty{-1.04e11}{kg}$. (b)~$1.04/1.74
= \qty{60}{\percent}$. Le reste se trouve dans les flancs de l'anomalie, hors
de la carte~: le granit atteint une profondeur de plus d'un kilomètre, et
l'anomalie de sa partie profonde décroît lentement avec la distance. Étendre la carte
assez loin, ou recourir à un modèle, permet de la restituer.
\end{solution}

\begin{solution}{exo:b3:field-lab-methods-3:7}
(a)~Pour un filon mince, l'anomalie tombe à la moitié de son pic en $x = \pm z$~:
sa largeur à mi-hauteur vaut $2z = \qty{60}{m}$, et trois
mesures en travers exigent un espacement d'au plus \qty{20}{m}. (b)~$2\times
220/3 = \qty{147}{m}$. (c)~L'espacement des stations requis dépend de la
profondeur du corps d'intérêt le moins profond~: une campagne conçue pour le
filon profond, à \qty{150}{m}, aliaserait l'anomalie du corps superficiel, qui
n'est large que de \qty{60}{m}~; une campagne conçue pour le corps superficiel voit aussi le
filon profond, avec bien plus de mesures qu'il n'en faut.
\end{solution}

\begin{solution}{exo:b3:field-lab-methods-3:8}
(a)~$F = 80/8 = 10$ et $\phi = F^{-1/m} = 10^{-1/2} = \num{0.316}$.
(b)~$10^{-1/1.8} = \num{0.278}$ et $10^{-1/2.2} = \num{0.351}$. Une variation
de \num{0.2} dans $m$ déplace la porosité d'environ quatre points de pourcentage~;
les résistivités, mesurées à quelques pour cent près, la déplacent bien moins. Le
exposant de cimentation contrôle l'incertitude, et il doit être calibré
sur des échantillons de la même roche.
\end{solution}

\begin{solution}{exo:b3:field-lab-methods-3:9}
(a)~Pour $s = 2h_1$, $(2nh_1/s)^2 = n^2$, les termes valent $k^n(1 +
n^2)^{-3/2}$~: $0.3536k$, $0.0894k^2$, $0.0316k^3$. Pour $k = 0.5$~:
$\num{0.1768}$, $\num{0.0224}$, $\num{0.0040}$, et $\rho_a/\rho_1 \approx
1 + 2\times0.2031 = \num{1.406}$ (la somme complète donne \num{1.409}). Pour $k
= -0.5$, les termes alternent en signe, $-0.1768 + 0.0224 - 0.0040$, et
$\rho_a/\rho_1 \approx \num{0.683}$ (somme complète \num{0.685}). (b)~Pour $s \ll
h_1$, tous les termes s'annulent et $\rho_a \to \rho_1$~; pour $s \gg h_1$, chaque
crochet tend vers~1 et $1 + 2\sum k^n = 1 + 2k/(1 - k) = (1 + k)/(1 -
k) = \rho_2/\rho_1$. (c)~$\rho_2/\rho_1 = (1 + k)/(1 - k) = 3$ et $1/3$.
\end{solution}

\begin{solution}{exo:b3:field-lab-methods-3:10}
(a)~Dans l'ordre AMNB avec un espacement $a$~: $AM = a$, $BM = 2a$, $AN = 2a$,
$BN = a$, donc le crochet vaut $2/a - 1/a = 1/a$ et $K = 2\pi a$.
(b)~Avec A et B en $\mp s$ et M et N en $\mp l$~: $AM = BN = s - l$,
$BM = AN = s + l$, le crochet vaut $2/(s - l) - 2/(s + l) = 4l/(s^2 -
l^2)$, et $K = \pi(s^2 - l^2)/(2l)$. Lorsque $l \ll s$, $\Delta V_{MN}
\approx 2l\,E$ avec $E = 2(-\mathrm{d}V/\mathrm{d}r)$ au centre, le
champ des deux électrodes de courant~; alors $\rho_a = K\Delta V/I \to
\pi s^2/(2l)\times 4l(-\mathrm{d}V/\mathrm{d}r)/I = 2\pi s^2(-
\mathrm{d}V/\mathrm{d}r)/I$. (c)~Le rapport entre le facteur approché et le
facteur exact est $s^2/(s^2 - l^2) = 1/(1 - 0.01) = 1.010$~: un pour cent.
\end{solution}

\begin{solution}{exo:b3:field-lab-methods-3:11}
(a)~Lecture~: \qty{0.010}{\milli\gal}. Hauteur~:
$0.1967\times0.05 = \qty{0.010}{\milli\gal}$. Masse volumique~: la correction de Bouguer
est $0.04191\,h\,d$ avec $d$ en \unit{g.cm^{-3}}, donc
$0.04191\times50\times0.1 = \qty{0.210}{\milli\gal}$. Terrain~:
\qty{0.200}{\milli\gal}. Total~: $\sqrt{0.010^2 + 0.010^2 + 0.210^2 +
0.200^2} = \qty{0.29}{\milli\gal}$. (b)~La masse volumique et la \omterm{def:b3:field-lab-methods-3:terrain}{correction
de terrain} dominent, chacune environ vingt fois l'erreur de lecture. Pour détecter
\qty{0.1}{\milli\gal}, la masse volumique de la dalle doit être mesurée
sur des échantillons ou ajustée aux données (la valeur qui reflète le moins la
topographie), et la \omterm{def:b3:field-lab-methods-3:terrain}{correction de terrain} calculée à partir d'un modèle numérique de terrain fin.
(c)~$\sqrt{0.010^2 + 0.010^2 + 0.210^2} = \qty{0.21}{\milli\gal}$.
\end{solution}

\begin{solution}{exo:b3:field-lab-methods-3:12}
(a)~La masse qui s'accorde à la lecture centrale est $\Delta m = \Delta g(0)z^2/G
= \num{-4.16e-5}z^2/\num{6.674e-11}$~: \qty{-5.61e10}{kg} à
\qty{300}{m}, \qty{-9.97e10}{kg} à \qty{400}{m}, \qty{-1.56e11}{kg} à
\qty{500}{m}. Les lectures prédites en $x = \pm300$ et $\pm600$ sont
\num{-1.47} et \num{-0.37}, \num{-2.13} et \num{-0.71}, \num{-2.62}
et \qty{-1.09}{\milli\gal}. Les écarts sont $\chi^2 = 1238$, $3.7$ et
$856$. (b)~Seul le modèle à \qty{400}{m}, avec $\chi^2 = 3.7$ pour cinq
données, s'inscrit dans les barres d'erreur. (c)~Oui~: tout corps à symétrie sphérique
de même masse et de même centre donne exactement la même anomalie, et d'autres
formes peuvent rester dans les erreurs. Une masse volumique mesurée sur des échantillons, un
forage ou une autre méthode trancherait.
\end{solution}

\begin{solution}{pb:b3:field-lab-methods-3:1}
\textbf{1.} $(2846.708 - 2846.512)/\qty{2.0}{h} = \qty{0.098}{\milli\gal.h^{-1}}$.

\textbf{2.} $2851.075 - 0.098\times0.5 - 2846.512 = \qty{4.514}{\milli\gal}$~;
$2850.497 - 0.088 - 2846.512 = \qty{3.897}{\milli\gal}$~;
$2849.392 - 0.127 - 2846.512 = \qty{2.753}{\milli\gal}$.

\textbf{3.} Les hauteurs au-dessus de la base sont $-27.1$, $-30.6$ et
\qty{-18.1}{m}~; en ajoutant $0.1967\,\Delta h$, on obtient les anomalies de Bouguer
$\num{-0.817}$, $\num{-2.122}$ et \qty{-0.808}{\milli\gal}. Les
stations sont plus basses que la base et plus proches du centre de la Terre,
donc leurs lectures brutes sont plus élevées~; la correction retire cet effet et
laisse un minimum au-dessus du granite.

\textbf{4.} L'écart $0.098 - 0.040 = \qty{0.058}{\milli\gal.h^{-1}}$,
soit \qty{0.12}{\milli\gal} sur la boucle, est la marée terrestre solide. Sur deux
heures, elle varie presque linéairement~; une droite passant par les deux
lectures à la base l'élimine avec la dérive.

\textbf{5.} Les erreurs sont $\num{0.008}$, $\num{0.009}$ et
\qty{0.017}{\milli\gal}. L'arrondi de chaque hauteur à \qty{0.1}{m} peut fausser
une différence de hauteur d'au plus \qty{0.1}{m}, soit
\qty{0.020}{\milli\gal}~: l'arrondi explique les erreurs, la marée
résiduelle après la droite étant inférieure à \qty{0.001}{\milli\gal}.

\textbf{6.} $0.01/0.1967 = \qty{5.1}{cm}$. Les \qty{3.3}{cm} des manuels valent
$0.01/0.3086$~: l'anomalie à l'air libre ne dépend pas de la masse volumique
choisie pour la correction de Bouguer, et la hauteur doit lui correspondre aussi.

\textbf{7.} $x_{1/2} = (3.952 - 3.048)/2 = \qty{0.452}{km}$, et $z = 1.305
\times452 = \qty{590}{m}$ sous le sol.

\textbf{8.} $\Delta m = \Delta g\,z^2/G = \num{-2.293e-5}\times590^2/
\num{6.674e-11} = \qty{-1.20e11}{kg}$. Avec $\Delta\rho =
\qty{-130}{kg.m^{-3}}$, son volume est \qty{9.2e8}{m^3} et son rayon
\qty{603}{m}, supérieur à la profondeur de son centre~: la sphère dépasserait
du sol. Aucune sphère avec le contraste de masse volumique du granite
ne convient~; le corps doit atteindre presque la surface.

\textbf{9.} Le granite s'étend de la surface jusqu'à une grande profondeur. Sa partie
superficielle, proche des stations, rend le pic étroit, et la règle, qui
ne connaît que la largeur, place une sphère unique à une profondeur modérée~; la
partie profonde ajoute des flancs larges que la sphère ne possède pas (la figure du
profil), et la masse qu'ils contiennent est manquée.

\textbf{10.} $2\pi G\Delta\rho = 2\pi\times\num{6.674e-11}\times(-130) =
\qty{-5.45e-8}{s^{-2}}$. Avec $z_1 = 0$, $z_2 = \qty{1645}{m}$, $a =
\qty{461}{m}$~: le crochet vaut $1645 + 461 - 1708.4 = \qty{397.6}{m}$,
et $\Delta g = \qty{-2.17}{\milli\gal}$. Pour une profondeur illimitée, le
crochet tend vers $a$, et $\Delta g = \qty{-2.51}{\milli\gal}$.

\textbf{11.} Le crochet doit valoir $2.29/\num{5.45e-3} = \qty{420}{m}$, ce qu'un
plancher vers \qty{2.6}{km} donne. Tout plancher entre environ
\qty{1.6}{km} et une profondeur illimitée ne modifie le pic que de quelques
dixièmes de milligal, et un corps légèrement plus large ou plus dense donne le même
pic avec un plancher moins profond~: le plancher n'est pas résolu.

\textbf{12.} Avec $z_1 = \qty{6}{m}$~: $1639 + 461.0 - 1708.4 = 391.6$, et
$\Delta g = \qty{-2.14}{\milli\gal}$, soit \qty{0.03}{\milli\gal} de moins que
avec le sommet au niveau du sol. C'est en dessous des erreurs de terrain et de
masse volumique de la campagne~: la gravimétrie ne peut pas dire si le granite atteint
la surface sous les alluvions. Le sondage le peut.

\textbf{13.} $K = \pi(20^2 - 1^2)/(2\times1) = \qty{626.7}{m}$, et $\Delta V =
\rho_a I/K = 821\times0.050/626.7 = \qty{65.5}{mV}$.

\textbf{14.} $V = 2\times1/\sqrt{17.1^2 - 10.7^2} = 2/13.34 =
\qty{0.150}{m.ns^{-1}}$, $\varepsilon_r = (0.2998/0.150)^2 = 4.0$, un
sable sec~; la canalisation est à $0.150\times10.7/2 = \qty{0.80}{m}$ de profondeur.

\textbf{15.} $0.150\times20.0/2 = \qty{1.50}{m}$.

\textbf{16.} $20/821 = \num{0.0244}$, $30/1309 = \num{0.0229}$, $40/1684 =
\num{0.0238}$~: $S \approx \qty{0.0237}{S}$.

\textbf{17.} La couche sèche contribue $1.5/600 = \qty{0.0025}{S}$, la
couche saturée $\qty{0.0212}{S}$, donc son épaisseur vaut $0.0212\times222 =
\qty{4.7}{m}$, et le granite se trouve $1.5 + 4.7 = \qty{6.2}{m}$ sous le
plancher. L'estimation est élevée~: le granite n'est pas infiniment
résistant, si bien qu'à ces écartements $\rho_a$ reste sous la droite
$s/S$, et $s/\rho_a$ surestime $S$. (Le modèle de la région comporte
\qty{6}{m} d'alluvions.)

\textbf{18.} $\phi = (20/222)^{1/2} = \num{0.30}$. Avec $m = 1.3$, $\phi =
0.0901^{1/1.3} = \num{0.16}$~; avec $m = 2.8$, $\phi = \num{0.42}$.

\textbf{19.} En retranchant la base de la ligne, on obtient, d'ouest en est~:
3.4, 4.1, 4.8, 5.9, 7.4, 9.2, 11.5, 14.2, 17.2, 19.8, 21.0, 20.1, 17.6,
14.5, 11.5, 9.1, 7.2, 5.7, 4.6, 3.8 et \qty{3.1}{nT}.

\textbf{20.} Sans correction, le pic vaut \qty{28.0}{nT} à \qty{2.20}{km}, et
sa moitié, \qty{14.0}{nT}, n'est jamais atteinte à l'est~: la ligne se termine à
\qty{15.2}{nT}, car le champ a augmenté de \qty{12}{nT} pendant la
marche.

\textbf{21.} Le pic vaut \qty{21.0}{nT} à \qty{2.20}{km}~; sa moitié vaut
\qty{10.5}{nT}. À l'ouest, entre \num{9.2} à \num{1.95} et \num{11.5} à
\qty{2.00}{km}, on a $1.95 + 0.05\times1.3/2.3 = \qty{1.978}{km}$. À l'est,
entre \num{11.5} à \num{2.40} et \num{9.1} à \qty{2.45}{km}, on a
$2.40 + 0.05\times1.0/2.4 = \qty{2.421}{km}$. La demi-largeur vaut $(2.421
- 1.978)/2 = \qty{0.221}{km}$.

\textbf{22.} Une anomalie symétrique dans un champ incliné à \ang{65} signifie que seule
l'aimantation verticale agit~: le corps se dirige selon le méridien magnétique,
nord--sud. La règle de la
\cref{prop:b3:field-lab-methods-3:dyke} suppose un dyke mince, vertical,
bidimensionnel et s'enfonçant loin, aimanté par induction, dont
le sommet agit comme une ligne de charge magnétique.

\textbf{23.} $\kappa t = 2\pi z\,\Delta T_\mathrm{max}/(F\sin^2 I) = 2\pi\times
221\times21.0/(48000\times0.8214) = \qty{0.74}{m}$. Avec $\kappa =
0.03$, $t = \qty{25}{m}$. Seul le produit est fixé~: un dyke deux fois
plus épais avec une susceptibilité moitié donne la même anomalie.

\textbf{24.} L'anomalie a \qty{443}{m} de large à mi-hauteur et les stations
sont espacées de \qty{50}{m}~: neuf lectures la couvrent, soit trois fois
le minimum de la règle.

\textbf{25.} Le sommet du dyke se trouve à
\textbf{\qty[text-series-to-math=true]{220}{m}} sous le sol. La
carte géologique place la discordance à \qty{222}{m} de profondeur en ce point~:
le dyke remonte jusqu'à la base de la couverture et s'y arrête, étant plus ancien
que le grès triasique qui l'enfouit.
\end{solution}
